\documentclass[11pt]{article}

%%METADATA
\title{GA1025 Macro Theory \\Solution to Problem Set 2}
\author{
Junbiao Chen\thanks{E-mail: jc14076@nyu.edu.}
}
\date{\today}


%%PACKAGES
\usepackage{mdframed} % For boxed environments
\usepackage{graphicx}
\usepackage{grffile}
\usepackage{tabularx}
\usepackage{setspace}
\usepackage{amsmath,amsthm,amssymb}
\usepackage[hyphens]{url}
\usepackage{natbib}
\usepackage[font=normalsize,labelfont=bf]{caption}
\usepackage[margin=1in]{geometry}
\usepackage{hyperref}
\hypersetup{colorlinks=true,urlcolor=blue,citecolor=blue}
\usepackage{stmaryrd}  %Package with \boxast command
\usepackage{enumerate}% http://ctan.org/pkg/enumerate %Supports lowercase Roman-letter enumeration
\usepackage{verbatim} %Package with \begin{comment} environment
%\usepackage{enumitem}
\usepackage{physics}
\usepackage{tikz}
\usepackage{listings}
\usepackage{upquote}
\usepackage{booktabs} %Package with \toprule and \bottomrule
\usepackage{etoc}     %Package with \localtableofcontents
\usepackage{placeins}    %Package that prevent repositioning the tables
\usepackage{multicol}
\usepackage{bm}
\usepackage{subfig}
\usepackage{csquotes}

\definecolor{dkgreen}{rgb}{0,0.6,0}
\definecolor{gray}{rgb}{0.5,0.5,0.5}
\definecolor{mauve}{rgb}{0.58,0,0.82}

\lstset{language=bash,
  frame=tb,
  aboveskip=3mm,
  belowskip=3mm,
  showstringspaces=false,
  columns=flexible,
  basicstyle={\small\ttfamily},
  numbers=none,
  numberstyle=\tiny\color{gray},
  keywordstyle=\color{blue},
  commentstyle=\color{dkgreen},
  stringstyle=\color{mauve},
  breaklines=true,
  breakatwhitespace=false,
  tabsize=3
}

\lstset{language=C,
  aboveskip=3mm,
  belowskip=3mm,
  showstringspaces=false,
  columns=flexible,
  basicstyle={\small\ttfamily},
  numbers=none,
  numberstyle=\tiny\color{gray},
  keywordstyle=\color{blue},
  commentstyle=\color{dkgreen},
  stringstyle=\color{mauve},
  breaklines=true,
  breakatwhitespace=false,
  tabsize=4
}

\definecolor{lightblue}{rgb}{0.68, 0.85, 0.9} %

%CUSTOM DEFINITIONS
\theoremstyle{definition}
\newtheorem{definition}{Definition}[section]
\newtheorem*{remark}{Remark}
\setcounter{secnumdepth}{3}
\usepackage{tikz}
\usetikzlibrary{arrows.meta}
\usetikzlibrary{automata, positioning, arrows, calc}

\tikzset{
	->,  % makes the edges directed
	>=stealth, % makes the arrow heads bold
	shorten >=2pt, shorten <=2pt, % shorten the arrow
	node distance=3cm, % specifies the minimum distance between two nodes. Change if n
	every state/.style={draw=blue!55,very thick,fill=blue!20}, % sets the properties for each ’state’ n
	initial text=$ $, % sets the text that appears on the start arrow
 }

%% PROPOSITION
% Define the Proposition environment
\newmdenv[
  innerleftmargin=10pt, 
  innerrightmargin=10pt,
  innertopmargin=10pt,
  innerbottommargin=10pt,
  linecolor=black, 
  linewidth=1pt,
  backgroundcolor=white, 
  roundcorner=5pt
]{propositionbox}

\newtheoremstyle{boldtitle} % Define a new theorem style
  {10pt} % Space above
  {10pt} % Space below
  {\itshape} % Body font
  {} % Indent amount
  {\bfseries} % Theorem head font
  {.} % Punctuation after theorem head
  { } % Space after theorem head
  {} % Theorem head spec

\theoremstyle{boldtitle} % Use the custom style
\newtheorem{proposition}{Proposition} % Define the proposition environment

% Redefine the proposition environment to use the box
\newenvironment{boxedproposition}[1][]
{\begin{propositionbox}\begin{proposition}[#1]}
{\end{proposition}\end{propositionbox}}

%%FORMATTING
\usepackage[bottom]{footmisc}
\onehalfspacing
\numberwithin{equation}{section}
\numberwithin{figure}{section}
\numberwithin{table}{section}
\bibliographystyle{../bib/aeanobold-oxford}


%main text
\begin{document}
\maketitle

\section{Pissarides’ analysis of taxation and variable search intensity}
\paragraph{a.} According to the setting, the timing protocol of the model is:
at the beginning of the period, the worker chooses the level of effort, 
then she instantaneously observes the probability of receiving an offer.
If the offer comes, she also needs to decide whether to accept the offer at the beginning of this period.
Therefore, the worker's problem can be characterized by the following value function.
\begin{align*}
    U = \max_{e} \left( 1 - \pi(e) \right) (1 - e + z + \beta U) + 
    \pi (e) \int_{0}^B \max_{\{\text{Accept, Reject}\}} \bigg\{& w - e - T(w) + \beta (\alpha U + (1 - \alpha) (w - T(w))), \\
    & 1 - e + z + \beta U \bigg\} dF(w)
\end{align*}
where $U$ denotes the value of unemployment.
%
%

The optimal strategy is given by the independent choices of optimal level and a reservation wages. 
In particular, 
at the reservation wage $\bar{w}$,
the indifference condition of accepting/rejecting an offer is 
\[
\bar{w} - e - T(\bar{w}) + \beta \bigg[\alpha U + (1 - \alpha)\left[\bar{w} - T(\bar{w}) \right] \bigg] = 1 - e + z + \beta U
\]
Rearranging terms, we have 
\[
\left[1 + \beta  (1 - \alpha)\right] \left[\bar{w} - T(\bar{w}) \right]  = 1 + z + \beta(1- \alpha) U
\]
Since $U$ is a constant, we can conclude that $\bar{w}$ does NOT depend on the optimal effort.

\vspace{8mm}
\noindent \textbf{Remark:} Intuitively, the choice of $e$ is a \textbf{sunk cost}, which won't affect the subsequent decisions.
%

To characterize the optimal choice of $e$, I first simplify the equation of $U$:
\begin{align*}
    U = \max_{e} 1 - e + z + \beta U - \pi(e) (1 + z + \beta U) + 
    \pi (e) \int_{0}^B \max_{\{\text{Accept, Reject}\}} \bigg\{& w - T(w) + \beta (\alpha U + (1 - \alpha) (w - T(w))), \\
    & 1 + z + \beta U \bigg\} dF(w)
\end{align*}

Next, I denote the objective function by $g(e)$, and denote 
$\int_{0}^B \max_{\{\text{Accept, Reject}\}} \bigg\{ w - T(w) + \beta (\alpha U + (1 - \alpha) (w - T(w))), 1 + z + \beta U \bigg\} dF(w)$ by $\Omega$.
I compute the derivative of $g(e)$ w.r.t. $e$ and obtain 
\[
\pi'(e) = \frac{1}{\Omega - (1 + z + \beta U)}
\]


\paragraph{b.}
Assume that $T(w) = t(w -a)$, then the indifference condition becomes 
\[
\left[1 + \beta  (1 - \alpha)\right] \left[(1-t) \bar{w} + ta \right]  = 1 + z + \beta(1- \alpha) U
\].
It follows that 
\[
\bar{w} = - \frac{t}{1 - t} a + \frac{1 + z + \beta (1 - \alpha )U}{1 + \beta (1 - \alpha) (1 - t)}
\].
Thus, we have $\frac{\partial \bar{w}}{\partial a} = - \frac{t}{1 - t} < 0$.

To examine the effect of increasing $a$ on the effort, 
I consider the function $h(e) = \pi'(e) - \frac{1}{\Omega - (1 + z + \beta U)}$.
Also note that $\frac{\partial \Omega}{\partial a} = \left( 1 - F(\bar{w})\right)(1 + \beta (1 - \alpha)t) > 0$.
Applying the implicit differenetial equation, I have 
\begin{align*}
\frac{\partial e^*}{\partial a} & = -\frac{\partial h/ \partial a}{\partial h/ \partial e^*} \\ 
& = - \frac{\frac{\frac{\partial \Omega}{\partial a}}{(\Omega - (1 + z + \beta U))^2}}{\pi''(e)} > 0
\end{align*}
This concludes that increasing $a$ increases the optimal level of search effort.
As the distribution of wage $F(w)$ is fixed, and the reservation wage decreases,
it immediately follows that the probability of accepting a job increases.

\paragraph{c.} 
The derivative of $\bar{w}$ with respect to $t$ is 
\[
\frac{\partial \bar{w}}{\partial t} = - \frac{a}{(1 - t)^2} + \frac{\beta (1 - \alpha)(1 + z + \beta (1-\alpha) U)
}{
    \left[1 + \beta (1 - \alpha)(1-t)^2 \right]^2
},
\]
which is positive if 
\[
\beta (1 - \alpha) \left[1 + z + \beta (1 - \alpha) U\right] > a \left[1 + \beta (1 - \alpha)\right]^2
\]

The derivative of $e^*$ with respect to $t$ is 
\begin{align*}
    \frac{\partial e^*}{\partial t} & = -\frac{\partial h/ \partial t}{\partial h/ \partial e^*} \\  
    & = - \frac{\beta (1 - \alpha)(a - \mathbb{E}_+[w]) + a - F(\bar{w}) [a + \beta a (1 - \alpha)] }{\pi''(e)}
\end{align*}
If the derivative is negative, we need 
\[
\frac{\beta (1 - \alpha)(a - \mathbb{E}_+[w]) + a - F(\bar{w}) [a + \beta a (1 - \alpha)] }{\pi''(e)} > 0
\]
That is, we need 
\[
\beta (1 - a) (a - \mathbb{E}_+[w]) + a < F(\bar{w})[a + \beta a (1 - \alpha)]
\]




\section{Jovanovic's model}%
\footnote{
    My code is available at \\ 
    \url{https://github.com/BillJunbiaoChen/GA1025_MacroTheory/tree/main/ProblemSets/PS2/PS2\_computation}.
}
\paragraph{a.} The value function is given by 
\begin{align*}
v(x_t) = \max_{s_t \geq 0, \phi_t \geq 0} \bigg\{
w_t(1 - s_t - \phi_t) + \beta \bigg[
& (1 - \pi(s_t)) v(x_{t+1}) + \\
& \pi (s_t) \int_0^B \max_{\text{Accept, Reject}} (v(\mu'), v(x_{t+1})) dF(\mu')    
\bigg]
\bigg\}
\end{align*}
\[
x_{t+1} = g(x_t \phi_t) - \delta x_t
\]


\paragraph{b.}
The decision rule to accept or reject an offer is given by a cutoff condition.
In particular, 
the worker will accept the offer if
\[
v(\mu) > v( g(x_t \phi_t) - \delta x_t)
\]
and will reject the offer if 
\[
v(\mu) < v( g(x_t \phi_t) - \delta x_t)
\]

\paragraph{c.}
I formulate the model using discretized $x_t$, $\phi$ and $s$.
Concretely, I formulate the quantitative model as follows:
\begin{table}[h!]
\centering
\caption{Parameterization and Variables of the Model}
\label{tab:parameterization}
\begin{tabular}{lll}
\hline
\textbf{Parameter}         & \textbf{Description}                            & \textbf{Value / Range}                     \\ \hline
$\beta$                    & Discount factor                                & 0.97                                       \\  
$\delta$                   & Depreciation rate                              & 0.1                                        \\  
$A$                        & Productivity of learning                       & 1.2                                        \\  
$\alpha$                   & Human capital investment elasticity            & 0.9                                        \\
$N$                        & Number of grid points                          & 100                                        \\
$f_i$           & Distribution weights                           & $\frac{1}{N}$                              \\
$\mu_{\text{max}}$         & Maximum human capital                          & 1.978                                      \\
$\phi_{\text{max}}$        & Maximum human capital investment               & 0.99                                       \\
$\mu_i$         & Grid for human capital                         & $[0.002, \mu_{\text{max}}]$ evenly spaced  \\
$\phi_i$        & Grid for human capital investment              & $[0.01, 1]$ with step size 0.02            \\
$s_i$           & Grid for search effort                         & $[0.01, 1]$ with step size 0.01            \\ \hline
\end{tabular}
\end{table}

For each current human capital $x_{t}$ and guess of the VF, 
I search for the optimal $\phi$ and $s$.
Note that $g(x_t \phi_t) - \delta x_t$ might generate values that falls in between discretized grids of $x$.
To resolve this ``overfloat'' problem, 
I use \texttt{Interpolations.linear\_interpolation} to approximate the value function at these points.
The results are presented in the figures below.

\begin{figure}[htbp]
  \centering
  \includegraphics[width=0.46\textwidth]{value_function.png}
  \includegraphics[width=0.46\textwidth]{serach_policy_function.png}
  \caption{Value function and search policy}
\end{figure}

\clearpage
\section*{2. Convergence results and contraction mapping}
\paragraph{a.} To show that the limit is unique, one can apply the Triangle Inequality. 

Concretely, for any positive $\epsilon$, there exists $N$, such that 
\[
\rho(x, y) \leq \rho(x_n, x) + \rho(x_n, y) < \frac{\epsilon}{2} + \frac{\epsilon}{2} = \epsilon
\]
Since $\epsilon$ can be arbitrary small, we conclude that $\rho(x, y) \rightarrow 0$.
Therefore, the limit is unique. \(\blacksquare \)

\paragraph{b.} To show that convergence implies Cauchy, one can apply the Triangle Inequality. 
Concretely, denote the limit by $x$, for any $m, k \in \mathbb{N}$, and any positive $\epsilon$, 
there exists an integer $N$, s.t.
\[
\rho(x_m, x_k) \leq \rho(x_m, x) + \rho(x_k, x) < \frac{\epsilon}{2} + \frac{\epsilon}{2} = \epsilon
\]
for any $m, k > N$.
Therefore, convergence implies Cauchy. \(\blacksquare \)


\paragraph{c.} Note that if $(x_n)$ is a Cauchy sequence,
then for any positive $\epsilon$, there exists $N$, such that 
$\rho(x_m, x_k) < \epsilon$ for any $m, k \geq N$.
Now take $k = N$,
we have $\rho(x_m, x_N) < \epsilon$ for any $m > N$.
By the inverse Triangle Inequality, we have 
\[
\rho(x_m, 0) < \rho(x_N, 0) + \epsilon
\]
Therefore, we bound the sequence by 
\[
\max\{ \max\{|x_1|, \ldots, |x_{N-1}|\}, |x_N| + \epsilon\}
\]

\paragraph{d.} We can prove the counter-positive statement, i.e., 
\[
\exists \text{ a subseq of } x_n \text{ does not converge to } x \Leftrightarrow x_n \text{does not converge to } x
\]
\paragraph{Proof}
``$\Rightarrow$''
Let $(x_{m_k})$ be a subsequence that does not converge to $x$.
This implies that, there exists a positive $\epsilon$, for any $M \in \mathbb{N}$, there exist some $m_k > M$, 
such that $|x_{m_k} - x| > \epsilon $. 
Since $m \rightarrow \infty$ as $m_k \rightarrow \infty$, we conclude that $x_n$ does not converge to $x$ as well. 

``$\Leftarrow$''
Let $x_n$ be a sequence that does not converge to $x$.
For a given $\epsilon > 0$,
there exists infinitely many $x_n$ such that $|x_n - x| > \epsilon$.
Therefore, one can construct a subsequence $(X_{m_k})$ by select elements from the set $\{x: |x_n - x| > \epsilon \}$.
\(\blacksquare\)



\subsection{A Pseudo-Cauchy Criterion}
For example, $S_m = \sum_{i=1}^m \frac{1}{i}$ is a Pseudo-Cauchy sequence, however, it does not converge.

\section*{3. Linear contraction mapping}
\paragraph{a.} ``Monotonicity'' requires that for $x1 > x2$, $g(x1) > g(x2)$,
which implies $b > 0$.
``Discounting'' requires that $\exists \beta \in (0, 1)$, such that,
$g(x+c) \leq g(x) + \beta c$, that is, 
\begin{align*}
& a + bx + bc \leq a + bx + \beta c \\
\Rightarrow & b \leq \beta
\end{align*}
In summary, $g(x)$ requires $b \in (0,1)$ to be a Contraction Mapping. 

\paragraph{b.} Since $g(x)$ will intersect $f(x) = x$ only once for any values of $a$ and $b$, 
we don't need to restrict $a$ and $b$ to guarantee the existance and uniqueness of the fixed point. 
This answer differs in that in \textbf{a.} because contraction is not a necessary and sufficient condition of getting a unique fixed point. 

\paragraph{c.} 
If $\mathcal{S} = [0, 1]$, the condition for a unqiue fixed point is given by 
\[
\frac{a}{1 - b} \in [0, 1]
\]
To show that the condition in \textbf{a.} won't guarantee a unique fixed piont if $\mathcal{S} = [0, 1]$,
we can consider cases where the fixed point is outside $[0, 1]$.




\section*{4. Contraction Mapping in the Search Model}
The Bellman Equation of the search model is given by 
\begin{align*}
(Tv)(w) = \max_{\text{Accept, Reject}} (\frac{w}{1 - \beta}, c + \beta \int_{0}^B v(w') dF(w'))
\end{align*}
\begin{itemize}
    \item ``Monotonicity'': Suppose $v_2 > v_1$, and denote the reservation wages by $\bar{w}_1$ and $\bar{w}_2$.
    It follows that 
    \begin{align*}
        & \int_{0}^{\bar{w}_1} \frac{\bar{w}_1}{1 - \beta} dF(w')  + \int_{\bar{w}_1}^B v_1(w') dF(w') \\
        < & \int_{0}^{\bar{w}_1} \frac{\bar{w}_1}{1 - \beta} dF(w')  + \int_{\bar{w}_1}^B v_2(w') dF(w') \\
        < & \int_{0}^{\bar{w}_2} \frac{\bar{w}_1}{1 - \beta} dF(w')  + \int_{\bar{w}_2}^B v_2(w') dF(w') 
    \end{align*}
    Therefore, it satisfies monotonicity.

    \item ``Discounting'': 
    \begin{align*} 
        (T(f+a))(w) & = \max_{\text{Accept, Reject}} (\frac{w}{1 - \beta}, c + \beta \int_{0}^B (f+a)(w') dF(w')) \\
        & = \max_{\text{Accept, Reject}} (\frac{w}{1 - \beta}, c + \beta \int_{0}^B f(w') dF(w')+ \beta a) \\
        & = \beta a + \max_{\text{Accept, Reject}} (\frac{w}{1 - \beta}, c + \beta \int_{0}^B f(w') dF(w')) \\
        & = \beta a + (Tf)(w)
    \end{align*}
\end{itemize}



\section*{5. Raising cattle}
\subsection*{Answer to Question 5.1}
Note that the law of motion for the herd size is given by $x_{t+1} = a (x_t - y_t), \forall t = 1,2, \ldots, (T-1)$ given $x_1 = h$.
Therefore, the sequence problem is given by 
\begin{align*}
& \max_{\{y_{t}\}_{t=1}^T} \bigg[\sum_{t=1}^T \beta^t \phi(y_t) \bigg] \\ 
\text{s.t.  } \quad & x_{t+1} = a(x_t - y_t) \\
& x_T = y_T
\end{align*}
where the state variable is $x_{t+1}$ and the control variable is $y_t$.
In fact, we can consolidate the intertemporal constraints and rewrite the model as follows:
\begin{align*}
& \max_{\{y_{t}\}_{t=1}^T} \bigg[\sum_{t=1}^T \beta^t \phi(y_t) \bigg] \\ 
& \text{s.t.  } \sum_{t=1}^T \quad \frac{y_t}{a^{t-1}} = h
\end{align*}


\subsection*{Answer to Question 5.2}
The Bellman equation is given by
\[
V_t(x) = \max_{x'} \phi(y) + \beta V_{t+1}(x')
\]


\subsection*{Answer to Question 5.3}
To examine the solution to this problem, I formulate the Lagragean as follows:
\[
\mathcal{L} = \bigg[\sum_{t=1}^T \beta^t \phi(y_t) \bigg] + \lambda \bigg(h - \sum_{t=1}^T \quad \frac{y_t}{a^{t-1}} \bigg)
\]
Then, the F.O.C for $y_t$ and $y_{t+1}$  implies that 
\[
\beta^t \phi'(y_t) = \lambda \frac{1}{a^{t-1}}
\]

\[
\beta^{t+1} \phi'(y_{t+1}) = \lambda \frac{1}{a^{t}}
\]

Combine these conditions together, we have 
\[
\frac{\beta \phi'(y_{t+1})}{\phi'(y_t)} = \frac{1}{a} \Rightarrow \phi'(y_t) = a \beta \phi'(y_{t+1})
\]

Applying the resource constraint iteratively, we have 
\begin{align*}
x_{2} & = ah - ay_1 \\
x_{3} & = a(x_2) - ay_2 = a^2 h - a^2 y_1 - a y_2 \\
\dots \\
x_t &= a^{t-1} h - \sum_{\tau = 1}^{t-1} {a^\tau y_{t-\tau}}
\end{align*}

Now we discuss the law of motion of $y_t$ by different combinations of $a$ and $\beta$.

If $a \beta < 1$, then $\phi'(y_{t+1}) > \phi'(y_t)$, and $y_{t+1} < y_t$. Therefore, $\phi(y)$ decreases overtime.
Under this condition, the size of cattles fall over times.

% Therefore, $\frac{x_{t+2}}{a} + y_{t+1} = a(x_t - y_t)$.
If $a \beta > 1$, then $\phi'(y_{t+1}) < \phi'(y_t)$, and $y_{t+1} > y_t$. Therefore, $\phi(y)$ increases overtime.
Under this condition, however, $x_{t}$ has two possible evolution paths.
The first one is that $x_{t}$ keeps decreasing till $T$.
The second one is that $x_{t}$ first increases and then decreases.
Note that, in both cases, $x_{t}$ is a concave function.


\subsection*{Answer to Question 5.4}
Consider a knife-edge case where the increase in $a$ changes $a \beta < 1$ to $a \beta > 1$ (with fixed $\beta$).
As such, $\Delta y_t = y_{t+1} - y_{t}$ increases.


\subsection*{Answer to Question 5.5}
Note that $\beta$ descreases as $r$ increases.
Consider a knife-edge case where a decrease in $\beta$ 
changes $a \beta > 1$ to $a \beta < 1$ (with fixed $a$).
As such, $\Delta y_t = y_{t+1} - y_{t}$ decreases.


\subsection*{Answer to Question 5.6}
If the revenue function is given by $\phi(y) = \phi_0 y^\alpha$, then the euler equation becomes 
\[
y_t = (a \beta)^{\frac{1}{\alpha - 1}} y_{t+1}
\]
Substitute all $y_{t}$ with $y_1$, we have $y_t = (a \beta)^{-\frac{t-1}{\alpha - 1}} y_1$.
Pluging them into the aggregate resource constraint, we have 
\begin{align*}
    & y_1 + \frac{(a \beta)^{-\frac{1}{\alpha -1}}}{a}y_1 + \ldots + \frac{(a \beta)^{-\frac{T-1}{\alpha -1}}}{a^{T-1}}y_1 = h \\
\Rightarrow \quad & y_1^* = h \left[ \frac{1 - a^{-1}(a \beta )^{-\frac{1}{\alpha - 1}}
    }{
        1 - a^{-T} (a \beta)^{-\frac{T}{\alpha - 1}}
    }\right]
\end{align*}
It follows that 
\[
y_t^* =  h \left[ \frac{(a \beta)^{-\frac{t-1}{\alpha - 1}} - a^{-1}(a \beta )^{-\frac{t}{\alpha - 1}} 
    }{
        1 - a^{-T} (a \beta)^{-\frac{T}{\alpha - 1}}
    }\right]
\]

Since $y_t^*$ does NOT depend on $\phi_0$, the $\phi_0$ won't affect 
the optimal sold cattles and stocks across times, but simply changes the level of sales.
Meanwhile, the effect of increasing $\alpha$ on $y_t^*$ depends on the value of $\alpha \beta$.
\begin{itemize}
    \item If $\alpha \beta = 1$, $\frac{\partial y_t^*}{\partial \alpha} = 0$ (namely, $y_t$ doesn't depend on $\alpha$.)
    \item If $\alpha \beta > 1$, $\frac{\partial y_t^*}{\partial \alpha} < 0$ for some early periods, and $\frac{\partial y_t^*}{\partial \alpha} > 0$ for some late periods.
    \item If $\alpha \beta < 1$, $\frac{\partial y_t^*}{\partial \alpha} > 0$ for some early periods, and $\frac{\partial y_t^*}{\partial \alpha} < 0$ for some late periods.
    Namley, \textbf{sales become front-loaded.}
\end{itemize}

\subsection*{Answer to Questions 5.7 and 5.8}
Note that 
\[
y_t^* =  h \left[ \frac{(a \beta)^{-\frac{t-1}{\alpha - 1}} - a^{-1}(a \beta )^{-\frac{t}{\alpha - 1}} 
    }{
        1 - (a^\alpha \beta)^{\frac{T}{1 - \alpha}}
    }\right]
\]
Therefore, 
\begin{itemize}
    \item If $a^\alpha \beta < 1$, then $y_t^* =  h \left[(a \beta)^{-\frac{t-1}{\alpha - 1}} - a^{-1}(a \beta )^{-\frac{t}{\alpha - 1}} \right]$ as $T \rightarrow \infty$.
    \item If $a^\alpha \beta > 1$, then $y_t^* =  0$ as $T \rightarrow \infty$.
\end{itemize}



%% =========== %% ========= %%
\bibliography{../bib/references.bib}

\end{document}